\section{분자와 분모의 벌 수를 맞춘다}

\claimx{앞 노트는 점추정을 스물다섯 벌로 내고 그 표준오차를 한 벌로 냈다 --- 벌 수를 맞추면 파괴 대조 여덟 칸의 표준오차 배수가 전부 커진다}{벌 수를 맞춘 표준오차가 한 벌 표준오차보다 좁아지지 않는 칸이 하나라도 있으면 이 주장은 떨어진다.}

집계 자를 정본으로 올리는 등록 규칙은 파괴 대조의 델타가 자기 표준오차의 두 배를 넘는가를 묻는다. 앞 노트의 델타는 겹 씨앗 다섯과 뒤섞기 씨앗 다섯을 곱한 \textbf{스물다섯 벌}의 평균인데, 그 표준오차를 내는 함수는 붓스트랩 복제마다 뒤섞기 씨앗을 \textbf{하나만} 썼다. 분자와 분모가 다른 것을 재고 있었다. \textbf{앞 노트 자신이 같은 사이클에서 「벌 수 없는 조건은 무효」를 등록 규칙으로 신설해 놓고 그 규칙을 어기는 조건을 목표 문서에 박았다.}

이 노트는 복제마다 같은 스물다섯 벌을 평균한 뒤 그 평균의 표준편차를 낸다. 여기서 자기 적발이 하나 나온다 --- 유보 재표집 함수가 겹 씨앗을 물기 때문에, 그대로 두면 한 복제가 유보 재표집 다섯 개의 평균이 되어 표준오차가 거짓으로 좁아진다. 정합 팔은 복제마다 유보 재표집을 하나로 고정했다. 벌 수를 맞추자 표준오차가 좁아진 칸이 $96 / 96$ 이다. 배선은 이 노트의 구판 팔이 앞 노트의 함수와 소수 여섯 자리까지 같은 수를 낸다는 것을 보인다 --- 배관을 안 바꾸고 벌 수만 바꿨다는 증거다.

\section{채택 문을 넓히고 선택 규칙에 검정력을 넣는다}

\claimx{앞 노트의 선택 규칙은 채택 문을 통과한 자 가운데 검정력이 가장 낮은 자를 골랐다 --- 그것은 그 앞 노트가 검정력 영짜리 자를 정본으로 올린 바로 그 원리다}{검정력으로 고른 자가 앞 노트가 고른 자와 같으면 이 주장은 떨어진다.}

앞 노트의 규칙은 채택 문을 학습 라벨 전량 파괴 두 칸으로만 두고, 통과한 자 가운데 \textbf{가장 큰 도메인의 가중 몫이 가장 작은 자}를 골랐다. 그런데 그 자가 여덟 칸에서 통과하는 칸은 \textbf{정확히 그 두 칸}이었고, 통과자 셋 가운데 검정력이 가장 낮았다. 이 노트는 규칙을 고친다: 여덟 칸을 전부 싣고, \textbf{검정력}(등록된 정칙화 두 값에 걸친 표준오차 배수의 최솟값)이 가장 큰 자를 고르며, 도메인 몫은 진단으로만 병기한다.

결과는 $R_iv SE² 역가중$ 이고 앞 노트가 고른 자는 $R_z 순열SE 역가중$ 였다 --- 정본 선택이 뒤집힌다. \textbf{다만 앞 노트가 그 앞 노트의 승격을 되돌린 판정은 살아남는다} --- 도메인 균등 가중은 벌 수를 맞춘 표준오차에서도 채택 문을 $1 / 2$ 통과한다. \textbf{그 대가는 숨기지 않는다} --- 검정력으로 고른 자는 유효 도메인 수가 $3.0455$ 이라 한 도메인이 절반을 넘게 진다. 이것은 멀티도메인 축에 대한 비용이고 다음 노트가 갚아야 할 값이다.

\section{자를 뽑기에서 뗀다 --- 그리고 자 넷은 셋이었다}

\claimx{도메인 안 순열 귀무의 표준편차는 닫힌 꼴이고, 그래서 귀무 분산 역수 가중은 행 가중과 같은 자다 --- 근사가 아니라 항등식이다}{뽑기 추정과 닫힌 꼴의 상대오차가 몬테카를로 잡음으로 설명되지 않으면 이 주장은 떨어진다.}

순위상관은 중간순위 위의 피어슨 상관이고, 두 고정 점수벡터를 무작위 순열로 짝지으면 그 상관의 분산은 \textbf{동률이 있든 없든 정확히} 표본 크기에서 하나 뺀 값의 역수다. 그러므로 도메인 안 순열 귀무의 표준편차는 그 값의 제곱근의 역수이고, 가중으로 쓰는 그 역수의 제곱은 \textbf{표본 크기에서 하나 뺀 값} 그 자체다. 뽑기 추정과 닫힌 꼴의 최대 상대오차는 $0.0271$ 이고 몬테카를로 오차의 이론값이 $0.01582$ 이라 그 $1.7138$ 배다 --- 어긋남이 전부 뽑기 잡음이다. 닫힌 꼴 가중과 행 가중의 가장 큰 도메인 몫 차는 $0.001899$ 이고 그 차는 통째로 \textbf{하나를 빼는 보정}에서 온다. \textbf{앞 노트가 산문으로 적은 「최소분산 가중은 행 가중으로 되돌아간다」는 항등식이었고, 네 자는 사실 표본 크기에서 하나 뺀 값을 지수로 올린 한 족의 세 점 --- 지수가 영, 이분의 일, 일인 셋 --- 이었다.}

\section{앞 노트의 등록물을 전부 채점한다}

\claimx{앞 노트가 게재한 크기 맞춤 대조는 강등된 자의 한 칸이고, 그 노트의 정본 자에서는 자기 반증조건에 걸린다}{정본 자의 두 정칙화 값에서 그 차가 짝 표준오차의 두 배를 넘으면 이 주장은 떨어진다.}

앞 노트는 예측 여덟을 등록하고 \textbf{둘만} 채점했다. 이 노트가 나머지를 채점한다. 크기 맞춤 대조 예측은 \textbf{어느 자로 재느냐에 매여 있다} --- 앞 노트가 정본으로 고른 자에서는 두 정칙화 값의 짝 표준오차 배수가 $1.6856$ 와 $2.262$ 로 \textbf{둘 다 문턱을 못 넘고}, 이 노트가 고른 자에서는 $2 / 2$ 다. \textbf{앞 노트는 그 수를 자기 정본 자가 아닌 자에서 뽑아 논문 주장 여섯에 실었고, 어느 자로 잰 수인지 적지 않았다 --- 여기서 그 주장을 자와 함께 다시 적는다.} 반증조건 쪽에서는 자 넷을 같이 내야 한다는 조건의 \textbf{분모가 결과를 본 뒤에 좁혀졌다}. 등록 문언대로 채점하면 $5 / 6$ 이고 그 노트가 신고한 값은 $5 / 5$ 였다. 분모를 좁힌 커밋은 측정이 돈 커밋보다 십사 분 사십일 초 뒤다.

\section{원천 쪽 --- 학습에 들어간 행에서 다시 잰다}

\claimx{바깥 원천이 안 쓰이는 까닭은 특징 거리가 아니라 혼합일 수 있다 --- 학습에 들어간 바깥 원천 행의 도메인 혼합이 유보 혼합과 사실상 무상관이다}{두 몫의 상관이 뚜렷한 양수이면 이 설명은 떨어진다.}

앞 노트는 표준화 평균차를 바깥 원천 \textbf{전량} $35,641$ 행에서 쟀다. 모형이 보는 것은 예산 배분이 고른 $1,710$ 행이다. 그 행에서 다시 재면 학습 혼합과 유보 혼합의 피어슨 상관이 $-0.1539$ 다. 정본 자 무게의 절반 이상을 지는 도메인이 학습의 $0.040936$ 이고, 학습을 지배하는 도메인은 유보의 $0.043663$ 다. 그리고 다섯 겹의 선택 행이 \textbf{바이트로 같다} --- 겹 씨앗의 변동은 안쪽 원천과 겹 배정에서만 온다. \textbf{앞 노트의 산출물이 이 수를 이미 다 들고 있었는데 안 셌다.} 덧붙여 앞 노트의 「특징 분포가 겹치지 않는다」는 과장이다 --- 두 정규분포로 보면 가장 안 겹치는 특징도 $0.7537$ 만큼 겹친다.
